\begin{textsample}{POS Dim 3 – ce – Score 21.00 – ce/ce\_inf\_15.txt}  \label{ex:f3_pos_072}
SÍNTESE DAS APRENDIZAGENS - Respeitar e expressar sentimentos e emoções. - Atuar em grupo e \textbf{demonstrar} interesse em construir novas relações, respeitando a diversidade e solidarizando -se com os outros. - Conhecer e respeitar regras de convívio social, manifestando respeito pelo outro. - Reconhecer a importância de ações e situações do cotidiano que contribuem para o \textbf{cuidado} de sua saúde e a manutenção de ambientes saudáveis. - Apresentar autonomia nas práticas de \textbf{higiene}, alimentação, vestir -se e no \textbf{cuidado} com seu bem-estar, valorizando o próprio corpo. - Utilizar o corpo intencionalmente ( com criatividade, controle e adequação ) como instrumento de interação com o outro e com o meio. - Coordenar suas habilidades manuais. - Discriminar os diferentes tipos de \textbf{sons} e ritmos e interagir com a música, percebendo -a como forma de expressão individual e coletiva. - Expressar -se por meio das artes visuais, utilizando diferentes materiais. - Relacionar -se com o outro empregando \textbf{gestos}, palavras, \textbf{brincadeiras}, jogos, \textbf{imitações}, observações e expressão corporal. - Expressar ideias, \textbf{desejos} e sentimentos em distintas situações de interação, por diferentes meios. - Argumentar e relatar fatos oralmente, em sequência temporal e causal, organizando e adequando sua fala ao contexto em que é produzida. - Ouvir, compreender, \textbf{contar}, recontar e criar narrativas. - Conhecer diferentes gêneros e portadores textuais, \textbf{demonstrando} compreensão da função social da escrita e reconhecendo a leitura como fonte de prazer e informação. - Identificar, nomear adequadamente e comparar as propriedades dos objetos, estabelecendo relações entre eles. - Interagir com o meio ambiente e com fenômenos naturais ou artificiais, \textbf{demonstrando} \textbf{curiosidade} e \textbf{cuidado} com relação a eles. - Utilizar vocabulário relativo às noções de grandeza ( maior, \textbf{menor}, igual etc. ), espaço ( dentro e fora ) e medidas ( comprido, curto, grosso, fino ) como meio de comunicação de suas experiências. - Utilizar unidades de medida ( dia e noite ; dias, semanas, \textbf{meses} e ano ) e noções de tempo ( presente, passado e futuro ; antes, agora e depois ), para responder a necessidades e questões do cotidiano. - Identificar e registrar quantidades por meio de diferentes formas de representação ( contagens, desenhos, símbolos, escrita de números, organização de gráficos básicos etc. ).
De acordo com a Síntese das Aprendizagens, a organização do currículo e da Proposta Pedagógica constitui -se um desafio para a \textbf{professora}, o professor, porque ele precisa pensar em práticas de articulação que produzam sentido para as \textbf{crianças}, com experiências que promovam as aprendizagens que precisam ser garantidas nesta etapa, de forma a não assumir caráter preparatório, mas de construção e concretização do desenvolvimento integral e integrado da \textbf{criança}.
Em cada campo de experiências, estão previstas aprendizagens para \textbf{bebês}, \textbf{crianças} bem \textbf{pequenas} e \textbf{crianças} \textbf{pequenas} e aspectos de seu desenvolvimento que servem de aspectos observáveis em suas experiências na Educação \textbf{Infantil}.
A fim de favorecer um processo de transição integrado e contínuo da Educação \textbf{Infantil} para o Ensino Fundamental, para isso, é fundamental construir uma documentação que possibilite à \textbf{criança} e sua família, bem como a equipe escolar que \textbf{acolherá} a \textbf{criança} na etapa seguinte, conhecer algumas vivências desse processo.
No campo O eu, o outro e o nós, estão sintetizadas grandes aprendizagens, tais como : - Respeitar e expressar sentimentos e emoções. - Atuar em grupo e \textbf{demonstrar} interesse em construir novas relações, respeitando a diversidade e solidarizando -se com os outros. - Conhecer e respeitar regras de convívio social, manifestando respeito pelo outro.
A resposta que cabe à Educação \textbf{Infantil} não é apenas dizer se a \textbf{criança} alcançou ou não esses objetivos, mas sim dizer : como a experiência vivida na Educação \textbf{Infantil} se concretizou, de forma a dar conta dessas aprendizagens. - Como as \textbf{crianças} e os \textbf{adultos} foram se preocupando, no decorrer do ano, em \textbf{demonstrar} respeito aos sentimentos e emoções de outras \textbf{crianças} e \textbf{adultos}? - Como foram sendo trabalhadas essas situações nas rodas de conversa? - E nas diferentes interações que as \textbf{crianças} foram estabelecendo no cotidiano? - Como as \textbf{crianças} foram sendo ajudadas a expressar sentimentos e emoções? - Como eram mediados os conflitos no espaço da unidade? - Como as \textbf{crianças} foram avançando na construção dessas aprendizagens que dizem respeito ao seu direito de “ aprender a conviver ”?
No quadro síntese das aprendizagens ( BNCC, Brasil, 2017 ), cada uma das aprendizagens nos remete a uma boa possibilidade de perguntas que nos ajuda a acompanhar, planejar e \textbf{documentar} os processos de aprendizagem e desenvolvimento das \textbf{crianças}.
É a clareza sobre como cada turma e como cada \textbf{criança} vivenciou as oportunidades promovidas na Educação \textbf{Infantil} que ajudará as/as professoras, os professores da etapa/ano/escola seguinte \textbf{acolher} melhor as \textbf{crianças} e dar continuidade à escolaridade.
A seguir estão listadas algumas estratégias que visam à articulação, integração curricular entre os anos da Educação \textbf{Infantil} e da Educação \textbf{Infantil} com o Ensino Fundamental ( família/creche, creche/pré-escola, pré-escola/anos iniciais do ensino fundamental ), a saber : - organização didática por projetos como elemento de continuidade, garantindo certas características comuns aos currículos de ambas as etapas, tais como : a escolha de atividades, de tempos e de espaços pelas \textbf{crianças} ; o trabalho em conjunto ; a importância dada aos conhecimentos e interesses de aprendizagem das \textbf{crianças}, bem como das representações que possuem do ensino fundamental ; o espírito investigativo presente na construção dos conhecimentos ; a integração das diversas experiências de aprendizagem ; - criação de processos, momentos de partilha e colaboração : reuniões de planejamento entre as professoras, os professores de anos diferentes ; reuniões com as famílias para acolhimento das suas expectativas em relação as mudanças para \textbf{creche}, \textbf{pré-escola} ou ensino fundamental e apresentar as informações sobre as propostas da nova etapa ; - organização de visitas das \textbf{crianças} em processo de transição à nova sala de atividades e/ou escola para realização de atividades em conjunto das \textbf{crianças} que se encontram em anos limites ; - discussão da documentação pedagógica, onde são registrados e dado visibilidade aos processos de aprendizagem e desenvolvimento das \textbf{crianças}, pela equipe pedagógica e professores da unidade escolar.
As estratégias apresentadas precisam ser tomadas como inspiração para realização, ampliação de outras ações que visam o bem-estar das \textbf{crianças} em processos de transição.
Para finalizar ressalta -se que a Educação \textbf{Infantil} e o Ensino Fundamental possuem finalidades diferenciadas que demarcam os seus territórios curriculares.
Mas alguns aspectos são semelhantes como o desenvolvimento integral da \textbf{criança} desde a Educação \textbf{Infantil} até o Ensino Fundamental.
Essas finalidades são essenciais para definição dos aspectos que as diferenciam como o delineamento dos aspectos comuns aos currículos dessas etapas e que garantem a continuidade e ampliação das experiências de aprendizagens iniciadas desde a Educação \textbf{Infantil} pelas \textbf{crianças} que prosseguem no Ensino Fundamental.

% matched lemmas: acolher, adulto, bebê, brincadeira, contar, creche, criança, cuidado, curiosidade, demonstrar, desejo, documentar, gesto, higiene, imitação, infantil, mês, pequeno, professora, pré-escola, som
\end{textsample}
